
% APF v35 -- Trace-to-Scheme Export Theorem
\section{Trace-to-Scheme Export: Codomain Admission, Transport, and Obstruction}
\label{sec:trace-to-scheme-export}

\paragraph{Purpose.}
The trace-sector evidence developed above shows that APF trace quantities are not, in general, physical observables.  They are locally closed sector anchors.  A physical comparison becomes meaningful only after a scheme/codomain, a transport map, and an uncertainty protocol have been supplied.  This section records that lesson as a general export theorem.

\begin{definition}[Trace anchor]
An APF trace anchor is a locally closed value or vector produced inside the APF trace algebra before assignment to an external physical measurement scheme.  Trace anchors may be dimensionless, such as $\sin^2\theta_W^{\rm APF\text{-}TRACE}$, or dimensionful, such as $M_W^{\rm APF\text{-}TRACE}$ and $m_f^{\rm APF\text{-}TRACE}$.
\end{definition}

\begin{definition}[Codomain admission]
Let $T$ be an APF trace anchor and let $S$ be an external scheme/codomain such as on-shell electroweak, self-scale $\overline{\rm MS}$, MSR, pole mass, Monte Carlo mass, or low-energy lattice $\overline{\rm MS}(2\,\mathrm{GeV})$.  The pair $(T,S)$ is codomain-admitted if the paper supplies a non-inverse, source-independent reason that $T$ may be interpreted as lying in, or being transported to, the codomain $S$.
\end{definition}

\begin{definition}[Admissible export route]
An admissible export route is a tuple
\[
\mathcal R=(T,S,\Phi,\Sigma,\mathcal A),
\]
where $T$ is the trace anchor, $S$ is the physical codomain, $\Phi:T\mapsto O_S$ is a declared transport map to an observable in $S$, $\Sigma$ is an uncertainty/covariance protocol, and $\mathcal A$ is a no-smuggling audit.  The route is admissible only if $\Phi$ does not consume the target observable being predicted or validated.
\end{definition}

\begin{definition}[Transport obstruction]
A transport obstruction is a named failed export gate: codomain absence, scheme ambiguity, missing normalization, missing finite-row evaluator, missing covariance, low-energy QCD/chiral obstruction, or target-observable smuggling.  A trace anchor may be locally closed while physically non-exported.
\end{definition}

\begin{lemma}[No-smuggling lemma]
A trace-to-scheme comparison is not evidence of APF export if the transport map, scale choice, normalization, or covariance protocol is fitted to the target observable without an independently declared route.  Equivalently, numerical agreement alone is insufficient; route admissibility is load-bearing.
\end{lemma}

\begin{theorem}[Trace-to-Scheme Export Theorem]
\label{thm:trace-to-scheme-export}
An APF trace anchor may be promoted from local trace closure to a physical export candidate only if all of the following gates are satisfied:
\begin{enumerate}
\item a codomain/scheme $S$ is declared;
\item a non-inverse transport map $\Phi$ is supplied or an identity codomain admission is justified;
\item all external constants and scale conventions used by $\Phi$ are listed;
\item uncertainty/covariance is propagated through the route;
\item the no-smuggling audit verifies that the target observable is not consumed as an input;
\item residuals are assigned to named transport, scheme, or covariance channels rather than hidden inside the trace value.
\end{enumerate}
If any gate fails, the trace anchor remains APF-meaningful at $[P_{\rm trace}]$ or $[P_{\rm local}]$, but the corresponding physical observable is not exported.  The failure is recorded as a transport obstruction rather than as a failure of trace closure.
\end{theorem}

\begin{proof}[Proof sketch]
A physical observable is not defined by a scalar value alone; it includes a measurement scheme, scale, renormalization convention, and uncertainty model.  Therefore equality or proximity between a trace anchor and an external number is not invariant under scheme changes.  By requiring codomain admission, a route map, constants ledger, covariance propagation, and a no-smuggling audit, the comparison becomes invariant under the declared route.  Conversely, if any of these ingredients is absent, the comparison is underdetermined: the same trace anchor can be routed to inequivalent physical codomains or fail to route at all.  Thus promotion to export-candidate status is equivalent to route completion, while failure of route completion yields a named obstruction.
\end{proof}

\paragraph{Case-study consequences.}
The theorem explains the observed sector pattern.  The electroweak $W$ branch closes as an export candidate because the route supplies a reviewed same-input DIZET transport evaluator, internal $\Delta r$ ledger, and covariance protocol.  The bottom branch closes as an export candidate because its trace anchor is admitted to the self-scale $\overline{\rm MS}$ codomain.  Charged leptons, top, charm, and light quarks remain locally trace-meaningful but physically non-exported or only externally codomain-admitted because their normalization, scheme, or low-energy QCD transport gates fail.

\begin{center}
\begin{tabular}{lll}
\hline
Sector & Closed status & Interpretation\\
\hline
EW/W & $[P_{\rm export\ candidate}]$ & reviewed same-input on-shell transport\\
$b$ & $[P_{\rm export\ candidate}]$ & self-scale $\overline{\rm MS}$ codomain admission\\
$e,\mu,\tau$ & $[P_{\rm residual\ obstruction}]$ & missing generation residual operator\\
$t$ & $[P_{\rm terminal\ scheme}]$ & external MSR lane only, no physical export\\
$c$ & $[P_{\overline{\rm MS}\ {\rm knockout}}]$ & missing charm normalization/transport map\\
$u,d,s$ & $[P_{\rm QCD\ obstruction}]$ & low-energy chiral/lattice transport required\\
\hline
\end{tabular}
\end{center}

\paragraph{Structural conclusion.}
APF predicts trace-sector anchors.  Physical observables require admissible route completion.  Thus a successful APF result need not be a direct scalar match; it may be an export candidate, a validation route, or a named obstruction that identifies the missing transport theorem.
